% === レイアウト・装飾 ===
% ---余白 ---
\geometry{
  top=25mm,
  bottom=25mm,
  left=20mm,
  right=20mm
}

% --- 句読点 ---
% \usepackage{newunicodechar}       % 「。、」を「．，」に変更
% \newunicodechar{、}{，}
% \newunicodechar{。}{．}

% --- インデント・段落間の空白 ---
% \setlength{\parindent}{0pt}       % インデントなし
% \setlength{\parskip}{0.6em}       % 段落間に空きを入れる

% --- プログラムコード ---
\lstset{
  basicstyle=\ttfamily\small,       % 基本のフォント
  commentstyle=\color{green!60!black},  % コメントの色
  keywordstyle=\color{blue},        % キーワードの色（int, forなど）
  stringstyle=\color{purple},       % 文字列の色
  frame=single,                     % 枠で囲む
  breaklines=true,                  % 長い行を折り返す
  numbers=left,                     % 行番号を左に
  numberstyle=\tiny\color{gray},    % 行番号のスタイル
  stepnumber=1,                     % 行番号を1行ごとに表示
  columns=fixed,                    % 文字幅を等間隔に
  basewidth=0.5em,
  showstringspaces=false,           % 文字列中のスペース記号を表示しない
}

% === 図表・式の番号設定 ===
\numberwithin{equation}{section}    % 数式番号を節ごとに
\counterwithin{table}{section}      % 表番号を節ごとに
\counterwithin{figure}{section}     % 図番号を節ごとに


% === 参考文献 ===
% --- thebibliography ----
% \usepackage[superscript]{cite}      % 参考文献

% --- BiBTeX ---
% \usepackage[backend=biber, style=numeric-comp, sorting=none]{biblatex}  % BiBLaTeX
% \addbibresource{references.bib}     % BiBLaTeXのファイル名